\scanpage{40}
On trouve ainsi un isom canonique
\[
\begin{aligned}
\Sigma^n(F/S)\simeq{}&\text{fibré affine associé à l'extension}\\
&0\longrightarrow\mathrm{Sym}^{n-1}(V)\otimes L\longrightarrow\mathrm{Sym}^n(V)\longrightarrow\mathcal O_S\longrightarrow0\\
&\text{(dont le fibré des translations}\\
&\text{est }\mathrm{Sym}^{n-1}(V)\otimes L\text{ )}
\end{aligned}\tag*{(62)}
\]
Cette suite exacte s'interprète comme celle
\inserted{\unsure{associé|associée} à l'hom} allant des formes de degré $n$ sur $\check V$ \inserted{$\simeq V\otimes L^{-1}$} vers $\mathcal O_S$,
par la ``valeur en $e_{\check V}=e_B$''. Donc la donnée
d'une section de $\Sigma^n(F/S)$, i.e. d'un
diviseur relatif de degré $n$ sur $F$, équivaut
à la donnée d'une forme de degré $n$ sur
$B=\check V$, prenant la valeur 1 en $e_B$.

6. Spécialisons ceci au cas de $n=2$. On trouve

\uline{Théorème} Soit $F$ fibré affine, associé à
ext. (47) de $\mathcal O_S$ par ``\deleted{fibré}\inserted{\unsure{module|espace}} des translations'' sur $L$,
ou dualement par ext. $B$ de $\check L$ par $\underline O_S$, correspondant
à la donnée, dans le Module loc. libre
\inserted{$B$} de rang 2, \deleted{$B$} d'une section $e_B$ non nulle
en chaque fibre. Il y a alors une
corr. biunivoque entre

a) Les structures d'Algèbre comm. unifère
sur $B$ admettant $e_B$ comme \unsure{section|Slection|élection|élément} unité.

b) Les diviseurs relatifs \inserted{\deleted{\unsure{$C$}} \unsure{$\Omega$|de} sur $F$} de degré 2 sur $S$

c) Les sections de $\Sigma^2(F/S)$
% archive sheet number: 39
